%% This document created by Scientific Word (R) Version 2.5
%% Starting shell: article


\documentclass[12pt,a4paper]{article}
\usepackage{amsmath}

%%%%%%%%%%%%%%%%%%%%%%%%%%%%%%%%%%%%%%%%%%%%%%%%%%%%%%%%%%%%%%%%%%%%%%%%%%%%%%%%%%%%%%%%%%%%%%%%%%%%%%%%%%%%%%%%%%%%%%%%%%%%%%%%%%%%%%%%%%%%%%%%%%%%%%%%%%%%%%%%%%%%%%%%%%%%%%%%%%%%%%%%%%%%%%%%%%%%%%%%%%%%%%%%%%%%%%%%%%%%%%%%%%%%
\usepackage{amssymb}
\usepackage{graphicx}
%\usepackage{sw20aip}

%TCIDATA{TCIstyle=article/art4.lat,lart,article}

%TCIDATA{OutputFilter=LATEX.DLL}
%TCIDATA{Version=5.50.0.2890}
%TCIDATA{<META NAME="SaveForMode" CONTENT="1">}
%TCIDATA{BibliographyScheme=Manual}
%TCIDATA{Created=Mon Feb 16 17:48:12 2004}
%TCIDATA{LastRevised=Thursday, June 14, 2007 12:43:24}
%TCIDATA{<META NAME="GraphicsSave" CONTENT="32">}

\include{DocumentDimensions}
\input{tcilatex}
\begin{document}

\author{Yves Peysson and Joan Decker}
\title{Flux surface averaging of plasma current and Ohmic electric field for CRONOS in LUKE}
\date{14th June 2007}
\maketitle

\subsection{\protect\bigskip Flux surface averaging definitions}

According to the LUKE notations (see LUKE documentation), the flux-surface averaging of a surface
quantity $\mathbf{\Gamma }\left( \psi ,\theta \right) $ along the toroidal
direction $\widehat{\phi }$ is 
\begin{equation}
\left\langle \Gamma \right\rangle _{\phi }\left( \psi \right) =\frac{1}{%
\overline{q}}\int_{0}^{2\pi }\frac{d\theta }{2\pi }\frac{1}{\left\vert 
\widehat{\psi }\cdot \widehat{r}\right\vert }\frac{r}{R}\frac{B_{0}}{B_{P}}%
\left[ \widehat{\phi }\mathbf{\cdot \Gamma }\left( \psi ,\theta \right) %
\right] 
\end{equation}%
while the flux-surface averaging of a volume quantity $\mathbf{\Phi }$ is 

\begin{equation}
\left\langle \mathbf{\Phi }\right\rangle _{V}\left( \psi \right) =\frac{1}{%
\widehat{q}}\int_{0}^{2\pi }\frac{d\theta }{2\pi }\frac{1}{\left\vert 
\widehat{\psi }\cdot \widehat{r}\right\vert }\frac{r}{R_{p}}\frac{B_{0}}{%
B_{P}}\mathbf{\Phi }\left( \psi ,\theta \right) 
\end{equation}

\subsection{Parallel Ohmic electric field}

By definition, 
\begin{align}
\left\langle E_{\Vert }\right\rangle _{\phi }\left( \psi \right) & =\frac{1}{%
\overline{q}}\int_{0}^{2\pi }\frac{d\theta }{2\pi }\frac{1}{\left\vert 
\widehat{\psi }\cdot \widehat{r}\right\vert }\frac{r}{R}\frac{B_{0}}{B_{P}}%
E_{\Vert }\left( \psi ,\theta \right) \left[ \widehat{\phi }\mathbf{\cdot }%
\widehat{b}\right]   \nonumber \\
& =\frac{1}{\overline{q}}\int_{0}^{2\pi }\frac{d\theta }{2\pi }\frac{1}{%
\left\vert \widehat{\psi }\cdot \widehat{r}\right\vert }\frac{r}{R}\frac{%
B_{0}}{B_{P}}\frac{B_{T}}{B}E_{\Vert }\left( \psi ,\theta \right)   \nonumber
\\
& =\frac{1}{\overline{q}}\int_{0}^{2\pi }\frac{d\theta }{2\pi }\frac{1}{%
\left\vert \widehat{\psi }\cdot \widehat{r}\right\vert }\frac{r}{R}\frac{%
B_{T}}{B_{P}}\frac{E_{\Vert }\left( \psi ,\theta \right) }{\Psi \left( \psi
,\theta \right) }
\end{align}%
Using relation 
\begin{equation}
E_{\Vert }\left( \psi ,\theta \right) =\frac{1}{\Psi \left( \psi ,\theta
\right) }\frac{R_{0}^{2}}{R^{2}}E_{\Vert 0}\left( \psi \right) 
\end{equation}%
where $E_{\Vert 0}\left( \psi \right) $ is the value at the minimum magnetic
field $B_{0}$, 
\begin{align}
\left\langle E_{\Vert }\right\rangle _{\phi }\left( \psi \right) & =\frac{1}{%
\overline{q}}\int_{0}^{2\pi }\frac{d\theta }{2\pi }\frac{1}{\left\vert 
\widehat{\psi }\cdot \widehat{r}\right\vert }\frac{r}{R}\frac{B_{T}}{B_{P}}%
\frac{E_{\Vert 0}\left( \psi \right) }{\Psi ^{2}\left( \psi ,\theta \right) }%
\frac{R_{0}^{2}}{R^{2}}  \nonumber \\
& =E_{\Vert 0}\left( \psi \right) \frac{1}{\overline{q}}\int_{0}^{2\pi }%
\frac{d\theta }{2\pi }\frac{1}{\left\vert \widehat{\psi }\cdot \widehat{r}%
\right\vert }\frac{r}{R}\frac{B_{T}}{B_{P}}\frac{1}{\Psi ^{2}\left( \psi
,\theta \right) }\frac{R_{0}^{2}}{R^{2}}
\end{align}%
or 
\begin{align}
\left\langle E_{\Vert }\right\rangle _{\phi }\left( \psi \right) & =E_{\Vert
0}\left( \psi \right) \frac{1}{\overline{q}}\frac{R_{p}}{R_{0}}%
\int_{0}^{2\pi }\frac{d\theta }{2\pi }\frac{1}{\left\vert \widehat{\psi }%
\cdot \widehat{r}\right\vert }\frac{r}{R_{p}}\frac{B}{B_{P}}\frac{\xi _{0}}{%
\xi }\left[ \frac{\xi }{\xi _{0}}\frac{B_{T}}{B}\frac{1}{\Psi ^{2}\left(
\psi ,\theta \right) }\frac{R_{0}^{3}}{R^{3}}\right]   \nonumber \\
& =E_{\Vert 0}\left( \psi \right) \frac{1}{\overline{q}}\frac{R_{p}}{R_{0}}%
\int_{0}^{2\pi }\frac{d\theta }{2\pi }\frac{1}{\left\vert \widehat{\psi }%
\cdot \widehat{r}\right\vert }\frac{r}{R_{p}}\frac{B}{B_{P}}\frac{\xi _{0}}{%
\xi }\times   \nonumber \\
& \left[ \frac{\xi }{\xi _{0}}\frac{B_{T}}{B_{T0}}\frac{B_{T0}}{B_{0}}\frac{%
B_{0}}{B}\frac{1}{\Psi ^{2}\left( \psi ,\theta \right) }\frac{R_{0}^{3}}{%
R^{3}}\right]   \nonumber \\
& =E_{\Vert 0}\left( \psi \right) \frac{1}{\overline{q}}\frac{B_{T0}}{B_{0}}%
\frac{R_{p}}{R_{0}}\int_{0}^{2\pi }\frac{d\theta }{2\pi }\frac{1}{\left\vert 
\widehat{\psi }\cdot \widehat{r}\right\vert }\frac{r}{R_{p}}\frac{B}{B_{P}}%
\frac{\xi _{0}}{\xi }\left[ \frac{\xi }{\xi _{0}}\frac{1}{\Psi ^{3}\left(
\psi ,\theta \right) }\frac{R_{0}^{4}}{R^{4}}\right]   \nonumber \\
& =E_{\Vert 0}\left( \psi \right) \frac{\widetilde{q}}{\overline{q}}\frac{%
R_{p}}{R_{0}}\frac{B_{T0}}{B_{0}}\lambda \sigma \left\{ \sigma \frac{\xi }{%
\xi _{0}}\frac{1}{\Psi ^{3}\left( \psi ,\theta \right) }\frac{R_{0}^{4}}{%
R^{4}}\right\}   \nonumber \\
& =E_{\Vert 0}\left( \psi \right) \frac{\widetilde{q}}{\overline{q}}\frac{%
R_{p}}{R_{0}}\frac{B_{T0}}{B_{0}}\overline{\lambda }_{1,-3,4}
\end{align}

In CRONOS code, only volume flux surface averaging is performed. Therefore%
\begin{equation}
\left\langle E_{\Vert }\right\rangle _{V}\left( \psi \right) =\frac{1}{%
\widehat{q}}\int_{0}^{2\pi }\frac{d\theta }{2\pi }\frac{1}{\left\vert 
\widehat{\psi }\cdot \widehat{r}\right\vert }\frac{r}{R_{p}}\frac{B_{0}}{%
B_{P}}E_{\Vert }\left( \psi ,\theta \right) 
\end{equation}

or%
\begin{equation}
\left\langle E_{\Vert }\right\rangle _{V}\left( \psi \right) =E_{\Vert
0}\left( \psi \right) \frac{1}{\widehat{q}}\int_{0}^{2\pi }\frac{d\theta }{%
2\pi }\frac{1}{\left\vert \widehat{\psi }\cdot \widehat{r}\right\vert }\frac{%
r}{R_{p}}\frac{B_{0}}{B_{P}}\frac{1}{\Psi \left( \psi ,\theta \right) }\frac{%
R_{0}^{2}}{R^{2}}
\end{equation}

It is then staightforward to introduce the bounce-averaged quantity 
\begin{equation}
\left\langle E_{\Vert }\right\rangle _{V}\left( \psi \right) =E_{\Vert
0}\left( \psi \right) \frac{1}{\widehat{q}}\int_{0}^{2\pi }\frac{d\theta }{%
2\pi }\frac{1}{\left\vert \widehat{\psi }\cdot \widehat{r}\right\vert }\frac{%
r}{R_{p}}\frac{B}{B_{P}}\frac{\xi _{0}}{\xi }\left[ \frac{\xi }{\xi _{0}}%
\frac{1}{\Psi ^{2}\left( \psi ,\theta \right) }\frac{R_{0}^{2}}{R^{2}}\right]
\end{equation}%
or%
\begin{equation}
\left\langle E_{\Vert }\right\rangle _{V}\left( \psi \right) =E_{\Vert
0}\left( \psi \right) \frac{\widetilde{q}}{\widehat{q}}\lambda \sigma
\left\{ \sigma \frac{\xi }{\xi _{0}}\frac{1}{\Psi ^{2}\left( \psi ,\theta
\right) }\frac{R_{0}^{2}}{R^{2}}\right\} 
\end{equation}

Then%
\begin{equation}
\left\langle E_{\Vert }\right\rangle _{V}\left( \psi \right) =E_{\Vert
0}\left( \psi \right) \frac{\widetilde{q}}{\widehat{q}}\overline{\lambda }%
_{1,-2,2}
\end{equation}%
and finally%
\begin{equation}
\frac{\left\langle E_{\Vert }\right\rangle _{V}\left( \psi \right) }{%
\left\langle E_{\Vert }\right\rangle _{\phi }\left( \psi \right) }=\frac{%
\overline{q}}{\widehat{q}}\frac{\overline{\lambda }_{1,-2,2}}{\overline{%
\lambda }_{1,-3,4}}\frac{R_{0}}{R_{p}}\frac{B_{0}}{B_{T0}}
\end{equation}

\paragraph{Case of circular concentric flux-surfaces}

In that case, since 
\begin{equation}
\overline{q}=\widehat{q}=\frac{\epsilon }{1+\epsilon }\frac{B}{B_{P}}
\end{equation}%
and%
\begin{equation}
\overline{\lambda }_{1,-3,4}=\overline{\lambda }_{1,-1,2}=\overline{\lambda }%
_{1,1,0}=\sqrt{\frac{1+\epsilon }{1-\epsilon }}
\end{equation}%
while%
\begin{equation}
\overline{\lambda }_{1,-2,2}=\overline{\lambda }_{1,0,0}=1
\end{equation}%
one finds%
\begin{eqnarray*}
\frac{\left\langle E_{\Vert }\right\rangle _{V}\left( \psi \right) }{%
\left\langle E_{\Vert }\right\rangle _{\phi }\left( \psi \right) } &=&\sqrt{%
\frac{1-\epsilon }{1+\epsilon }}\frac{R_{0}}{R_{p}}\frac{B_{0}}{B_{T0}} \\
&=&\sqrt{1-\epsilon ^{2}}\frac{B_{0}}{B_{T0}} \\
&\simeq &\frac{B_{0}}{B_{T0}}
\end{eqnarray*}%
since 
\begin{equation}
\frac{R_{0}}{R_{p}}=1+\epsilon 
\end{equation}

\subsection{Current Density}

The density of current carried by electrons is given by 
\begin{equation}
\mathbf{J}\left( \mathbf{x}\right) =q_{e}\iiint d^{3}p\;\mathbf{v}f\left( 
\mathbf{x},\mathbf{p}\right)
\end{equation}
so that the parallel current density is 
\begin{equation}
J_{\Vert}\left( \mathbf{x}\right) =q_{e}\iiint d^{3}p\;v_{\Vert}f\left( 
\mathbf{x},\mathbf{p}\right)
\end{equation}
which becomes in $\left( \psi,\theta,p,\xi\right) $ phase space 
\begin{equation}
J_{\Vert}\left( \psi,\theta\right) =2\pi
q_{e}\int_{0}^{\infty}p^{2}dp\int_{-1}^{1}d\xi\;\frac{p\xi}{\gamma m}f\left(
\psi,\theta,p,\xi\right)
\end{equation}

\subsubsection{Flux-Surface Averaging}

We are usually interested in the flux-surface averaged current density in
the toroidal direction. It is given by  
\begin{align}
\left\langle J_{\Vert}\right\rangle _{\phi}\left( \psi\right) & =\frac {1}{%
\overline{q}}\int_{0}^{2\pi}\frac{d\theta}{2\pi}\frac{1}{\left| \widehat{\psi%
}\cdot\widehat{r}\right| }\frac{r}{R}\frac{B_{0}}{B_{P}}J_{\Vert}\left(
\psi,\theta\right) \left[ \widehat{\phi}\mathbf{\cdot }\widehat{b}\right] 
\nonumber \\
& =\frac{1}{\overline{q}}\int_{0}^{2\pi}\frac{d\theta}{2\pi}\frac{1}{\left| 
\widehat{\psi}\cdot\widehat{r}\right| }\frac{r}{R}\frac{B_{0}}{B_{P}}
J_{\Vert}\left( \psi,\theta\right) \frac{B_{T}}{B}
\end{align}
and finally, using
\begin{equation}
\left\langle J_{\Vert}\right\rangle _{\phi}\left( \psi\right) =\frac {1}{%
\overline{q}}\int_{0}^{2\pi}\frac{d\theta}{2\pi}\frac{1}{\left| \widehat{\psi%
}\cdot\widehat{r}\right| }\frac{r}{R}\frac{B_{T}}{B_{P}}\frac{
J_{\Vert}\left( \psi,\theta\right) }{\Psi\left( \psi,\theta\right) }
\end{equation}

\subsubsection{Fokker-Planck Equation}

When we consider only the zero order distribution function, we have that $%
f_{0}$ is constant along a field line, so that 
\begin{equation}
f_{0}\left( \psi,\theta,p,\xi\right) =f_{0}^{(0)}\left( \psi,p,\xi
_{0}\right)
\end{equation}
where 
\begin{equation}
\xi_{0}=\sigma\sqrt{1-\frac{1}{\Psi\left( \psi,\theta\right) }\left(
1-\xi^{2}\right) }  \label{xi0}
\end{equation}

Consequently, we find 
\begin{align}
J_{\Vert}^{0}\left( \psi,\theta\right) & =2\pi
q_{e}\int_{0}^{\infty}p^{2}dp\int_{-1}^{1}d\xi\;\frac{p\xi}{\gamma m_{e}}
f_{0}\left( \psi ,\theta,p,\xi\right)  \nonumber \\
& =2\pi q_{e}\int_{0}^{\infty}p^{2}dp\int_{-1}^{1}d\xi\;\frac{p\xi}{\gamma
m_{e}}f_{0}^{(0)}\left( \psi,p,\xi_{0}\right)  \nonumber \\
& =2\pi q_{e}\int_{0}^{\infty}p^{2}dp\int_{-1}^{1}d\xi_{0}\Psi\left(
\psi,\theta\right) \times  \nonumber \\
& H\left( \left| \xi_{0}\right| -\sqrt{1-\frac{1}{\Psi\left( \psi
,\theta\right) }}\right) \;\frac{p\xi_{0}}{\gamma m_{e}}f_{0}^{(0)}\left(
\psi,p,\xi_{0}\right)  \label{jpar}
\end{align}
where the condition 
\begin{equation}
\left| \xi_{0}\right| \geq\sqrt{1-\frac{1}{\Psi\left( \psi,\theta\right) }}
\label{xi0cond}
\end{equation}
results from the equation (\ref{xi0}) and means that only the particle who
reach the position $\theta$ must be considered. Note that the integrand in
the equation (\ref{jpar}) is odd in $\xi_{0}$ for trapped electrons, since $%
f_{0}^{(0)}$ is symmetric in the trapped region. As a consequence, the
contribution from trapped electrons vanishes, and (\ref{jpar}) can be
rewritten as 
\begin{equation}
J_{\Vert}^{0}\left( \psi,\theta\right) =2\pi
q_{e}\int_{0}^{\infty}p^{2}dp\int_{-1}^{1}d\xi_{0}\Psi\left(
\psi,\theta\right) H\left( \left| \xi_{0}\right| -\xi_{0T}\right) \;\frac{
p\xi_{0}}{\gamma m_{e}}f_{0}^{(0)}\left( \psi,p,\xi_{0}\right)
\end{equation}

Therefore, the flux-surface averaged current density 
\begin{equation}
\left\langle J_{\Vert}\right\rangle _{\phi}^{0}\left( \psi\right) =\frac {1}{%
\overline{q}}\int_{0}^{2\pi}\frac{d\theta}{2\pi}\frac{1}{\left| \widehat{\psi%
}\cdot\widehat{r}\right| }\frac{r}{R}\frac{B_{T}}{B_{P}}\frac{
J_{\Vert}^{0}\left( \psi,\theta\right) }{\Psi\left( \psi ,\theta\right) }
\end{equation}
becomes 
\begin{align}
\left\langle J_{\Vert}\right\rangle _{\phi}^{0}\left( \psi\right) & =\frac{
2\pi q_{e}}{m_{e}}\int_{0}^{\infty}dp\;\frac{p^{3}}{\gamma}\frac {1}{%
\overline{q}}\int_{0}^{2\pi}\frac{d\theta}{2\pi}\frac{1}{\left| \widehat{\psi%
}\cdot\widehat{r}\right| }\frac{r}{R}\frac{B_{T}}{B_{P}}\frac {1}{\Psi\left(
\psi,\theta\right) }\times  \nonumber \\
& \int_{-1}^{1}d\xi_{0}\Psi\left( \psi,\theta\right) H\left( \left|
\xi_{0}\right| -\xi_{0T}\right) \;\xi_{0}f_{0}^{(0)}\left( \psi,p,\xi
_{0}\right)
\end{align}

The integrals over $\theta $ and $\xi _{0}$ can be permuted 
\begin{align}
\left\langle J_{\Vert }\right\rangle _{\phi }^{0}\left( \psi \right) & =%
\frac{2\pi q_{e}}{m_{e}}\int_{0}^{\infty }dp\;\frac{p^{3}}{\gamma }%
\int_{-1}^{1}d\xi _{0}\;H\left( \left\vert \xi _{0}\right\vert -\xi
_{0T}\right) \;\xi _{0}f_{0}^{(0)}\left( \psi ,p,\xi _{0}\right) \times  
\nonumber \\
& \frac{1}{\overline{q}}\int_{0}^{2\pi }\frac{d\theta }{2\pi }\frac{1}{%
\left\vert \widehat{\psi }\cdot \widehat{r}\right\vert }\frac{r}{R}\frac{%
B_{T}}{B_{P}}
\end{align}

We recognize the expression of the safety factor $q$ so that 
\begin{equation}
\left\langle J_{\Vert }\right\rangle _{\phi }^{0}\left( \psi \right) =\frac{%
2\pi q_{e}}{m_{e}}\frac{q}{\overline{q}}\int_{0}^{\infty }dp\;\frac{p^{3}}{%
\gamma }\int_{-1}^{1}d\xi _{0}\;H\left( \left\vert \xi _{0}\right\vert -\xi
_{0T}\right) \;\xi _{0}f_{0}^{(0)}\left( \psi ,p,\xi _{0}\right) \;
\end{equation}

Concerning the CRONOS approach of the flux surface averaging

\begin{equation}
\left\langle J_{\Vert }\right\rangle _{V}^{0}\left( \psi \right) =\frac{1}{%
\widehat{q}}\int_{0}^{2\pi }\frac{d\theta }{2\pi }\frac{1}{\left\vert 
\widehat{\psi }\cdot \widehat{r}\right\vert }\frac{r}{R_{p}}\frac{B_{0}}{%
B_{P}}J_{\Vert }^{0}\left( \psi ,\theta \right) 
\end{equation}%
which becomes%
\begin{eqnarray}
\left\langle J_{\Vert }\right\rangle _{V}^{0}\left( \psi \right)  &=&\frac{%
2\pi q_{e}}{m_{e}}\int_{0}^{\infty }dp\;\frac{p^{3}}{\gamma }\frac{1}{%
\widehat{q}}\int_{0}^{2\pi }\frac{d\theta }{2\pi }\frac{1}{\left\vert 
\widehat{\psi }\cdot \widehat{r}\right\vert }\frac{r}{R_{p}}\frac{B_{0}}{%
B_{P}}\times   \nonumber \\
&&\int_{-1}^{1}d\xi _{0}\;\Psi \left( \psi ,\theta \right) H\left(
\left\vert \xi _{0}\right\vert -\xi _{0T}\right) \;\xi _{0}f_{0}^{(0)}\left(
\psi ,p,\xi _{0}\right) 
\end{eqnarray}%
Permuting the integrals over $\theta $ and $\xi _{0}$ which are independent 
\begin{eqnarray}
\left\langle J_{\Vert }\right\rangle _{V}^{0}\left( \psi \right)  &=&\frac{%
2\pi q_{e}}{m_{e}}\int_{0}^{\infty }dp\;\frac{p^{3}}{\gamma }%
\int_{-1}^{1}d\xi _{0}H\left( \left\vert \xi _{0}\right\vert -\xi
_{0T}\right) \;\xi _{0}f_{0}^{(0)}\left( \psi ,p,\xi _{0}\right) \times  
\nonumber \\
&&\;\frac{1}{\widehat{q}}\int_{0}^{2\pi }\frac{d\theta }{2\pi }\frac{1}{%
\left\vert \widehat{\psi }\cdot \widehat{r}\right\vert }\frac{r}{R_{p}}\frac{%
B}{B_{P}}
\end{eqnarray}%
or%
\begin{equation}
\left\langle J_{\Vert }\right\rangle _{V}^{0}\left( \psi \right) =\frac{2\pi
q_{e}}{m_{e}}\frac{\widetilde{q}}{\widehat{q}}\int_{0}^{\infty }dp\;\frac{%
p^{3}}{\gamma }\int_{-1}^{1}d\xi _{0}H\left( \left\vert \xi _{0}\right\vert
-\xi _{0T}\right) \;\xi _{0}f_{0}^{(0)}\left( \psi ,p,\xi _{0}\right) 
\end{equation}

Consequently%
\begin{equation}
\frac{\left\langle J_{\Vert }\right\rangle _{V}\left( \psi \right) }{%
\left\langle J_{\Vert }\right\rangle _{\phi }\left( \psi \right) }=\frac{%
\widetilde{q}}{\widehat{q}}\frac{\overline{q}}{q}
\end{equation}

\paragraph{Case of circular concentric flux-surfaces}

In that case, the safety factor expression becomes 
\begin{equation}
q\left( r\right) =\frac{\epsilon}{\sqrt{1-\epsilon^{2}}}\frac{B_{T}}{B_{P}}
\label{qcirc2}
\end{equation}
with $\epsilon=r/R_{p}$ the inverse aspect ratio.

In addition, $\overline{q}\left( r\right) $ becomes 
\begin{align}
\overline{q}\left( r\right) & =\int_{0}^{2\pi}\frac{d\theta}{2\pi}\frac {r}{%
R_{p}}\frac{B_{0}}{B_{P}}  \nonumber \\
& =\int_{0}^{2\pi}\frac{d\theta}{2\pi}\frac{r}{R_{p}}\frac{B_{0}}{B}\frac {B%
}{B_{P}}  \nonumber \\
& =\epsilon\frac{B}{B_{P}}\frac{R_{p}}{R_{0}}  \nonumber \\
& =\frac{\epsilon}{1+\epsilon}\frac{B}{B_{P}}
\end{align}
since $R=R_{p}+r\cos\theta,$ and $B_{0}/B=R/R_{0}.$ We have then 
\begin{equation}
\frac{q\left( r\right) }{\overline{q}\left( r\right) }=\sqrt {\frac{
1+\epsilon}{1-\epsilon}}\frac{B_{T}}{B}
\end{equation}

In the case when $B_{T}\gg B_{P}$, we retrieve the bounce-averaged
coefficient $s^{\ast }$ and in the large aspect ratio limit $\epsilon \ll 1,$%
\begin{equation}
\lim_{\epsilon \rightarrow 0}\frac{q\left( r\right) }{\overline{q}\left(
r\right) }=\left( 1+\epsilon \right) \frac{B_{T}}{B}
\end{equation}

Consequently,%
\begin{eqnarray}
\frac{\left\langle J_{\Vert }\right\rangle _{V}\left( \psi \right) }{%
\left\langle J_{\Vert }\right\rangle _{\phi }\left( \psi \right) } &=&\frac{%
\widetilde{q}}{\widehat{q}}\sqrt{\frac{1-\epsilon }{1+\epsilon }}\frac{B}{%
B_{T}}  \nonumber \\
&=&\sqrt{1-\epsilon ^{2}}\frac{B}{B_{T}}  \nonumber \\
&\simeq &\frac{B}{B_{T}}  \nonumber \\
&=&\frac{B_{0}}{B_{T0}}
\end{eqnarray}%
as for the Ohmic electric field, since $\widetilde{q}=\widehat{q}\left(
1+\epsilon \right) .$

\end{document}
